\documentclass[12pt]{article}
\usepackage[a4paper,margin=1in]{geometry}
\usepackage[utf8]{inputenc}
\usepackage[T1]{fontenc}
\usepackage{amsmath,amssymb}
\usepackage{textcomp}
\usepackage{parskip}
\begin{document}
\begin{center}
{\Large \textbf{2015 STEP II}}
\end{center}

\subsection*{Page 1}
\begin{flushleft}
\bigskip\textbf{Section A: Pure Mathematics}\par
\bigskip\textbf{Question 1.} (i) By use of calculus, show that x - ln(1 + x) is positive for all positive x. Use this result\par
to show that\par
n\par
\[> ln(n + 1) .\]
k\par
\[k=1\]
\[\leq\]
\textbf{(ii)} By considering x + ln(1 - x), show that\par
\infty\par
\[< 1 + ln 2 .\]
\[k^2\]
\[k=1\]
\[\leq\]
\bigskip\textbf{Question 2.} In the triangle ABC, angle BAC = α and angle CBA = 2α, where 2α is acute, and BC = x.\par
\[Show that AB = (3 - 4 sin α)x.\]
The point D is the midpoint of AB and the point E is the foot of the perpendicular from C\par
to AB. Find an expression for DE in terms of x.\par
The point F lies on the perpendicular bisector of AB and is a distance x from C. The points\par
F and B lie on the same side of the line through A and C. Show that the line F C trisects\par
the angle ACB.\par
\bigskip\textbf{Question 3.} Three rods have lengths a, b and c, where a < b < c. The three rods can be made into a\par
\[triangle (possibly of zero area) if a + b \leq c.\]
Let T be the number of triangles that can be made with three rods chosen from n rods of\par
n\par
lengths 1, 2, 3, . . . , n (where n \leq 3). Show that T8 - T7 = 2 + 4 + 6 and evaluate T8 - T6.\par
\[Write down expressions for T2m - T2m-1 and T2m - T2m-2.\]
\[Prove by induction that T2m =\]
m(m - 1)(4m + 1) , and find the corresponding result for\par
an odd number of rods.\par
\end{flushleft}
\subsection*{Page 2}
\begin{flushleft}
\bigskip\textbf{Question 4.} (i) The continuous function f is defined by\par
\[tan f(x) = x (-\infty < x < \infty)\]
\[and f(0) = \pi. Sketch the curve y = f(x) .\]
\textbf{(ii)} The continuous function g is defined by\par
x\par
\[tan g(x) = (-\infty < x < \infty)\]
\[1 + x^2\]
x\par
\[and g(0) = \pi. Sketch the curves y = and y = g(x) .\]
\[1 + x^2\]
\textbf{(iii)} The continuous function h is defined by h(0) = \pi and\par
x\par
\[tan h(x) = (x \ne ±1) .\]
\[1 - x^2\]
\[(The values of h(x) at x = ±1 are such that h(x) is continuous at these points.) Sketch\]
x\par
\[the curves y = and y = h(x).\]
\[1 - x^2\]
\bigskip\textbf{Question 5.} In this question, the arctan function satisfies 0 \geq arctan x < 1 \pi for x \leq 0 .\par
\textbf{(i)} Let\par
n\par
\[S = arctan ,\]
\[n 2m2\]
\[m=1 \geq (\]
\[\leq\]
\[for n = 1, 2, 3, . . . . Prove by induction that\]
n\par
\[tan S = .\]
n\par
n + 1\par
Prove also that\par
n\par
\[S = arctan .\]
n\par
n + 1\par
2 4\par
\textbf{(ii)} In a triangle ABC, the lengths of the sides AB and BC are 4n and 4n -1, respectively,\par
\[and the angle at B is a right angle. Let angle BCA = 2α . Show that\]
n\par
\infty\par
\[α = \pi .\]
n 4\par
\[n=1\]
\[\leq\]
\end{flushleft}
\subsection*{Page 3}
\begin{flushleft}
\bigskip\textbf{Question 6.} (i) Show that\par
2 1 1\par
\[sec \pi - x = .\]
4 2\par
1 + sin x\par
\[1 ) \int\]
Hence integrate with respect to x.\par
1 + sin x\par
\textbf{(ii)} By means of the substitution y = \pi - x, show that\par
\[\pi \pi \pi\]
\[xf(sin x) dx = f(sin x) dx,\]
0 2 0\par
\[\sum \sum\]
where f is any function for which these integrals exist.\par
Hence evaluate\par
\[\pi x\]
dx .\par
0 1 + sin x\par
\[\sum\]
\textbf{(iii)} Evaluate\par
\[\pi 2x 3 - 3\pix 2\]
dx.\par
0 (1 + sin x)2\par
\[\sum\]
\bigskip\textbf{Question 7.} A circle C is said to be bisected by a curve X if X meets C in exactly two points and these\par
points are diametrically opposite each other on C.\par
\textbf{(i)} Let C be the circle of radius a in the x-y plane with centre at the origin. Show, by\par
giving its equation, that it is possible to find a circle of given radius r that bisects C\par
\[provided r > a. Show that no circle of radius r bisects C if r \geq a .\]
\textbf{(ii)} Let C1 and C2 be circles with centres at (-d, 0) and (d, 0) and radii a1 and a2, respec-\par
\[tively, where d > a^1 and d > a^2. Let D be a circle of radius r that bisects both C^1\]
2 2\par
a - a\par
2 1\par
\[and C^2. Show that the x-coordinate of the centre of D is .\]
4d\par
\[Obtain an expression in terms of d, r, a^1 and a^2 for the y-coordinate of the centre\]
of D, and deduce that r must satisfy\par
\[16r 2 d 2 \leq 4d 2 + (a^2 - a^1) 2 4d 2 + (a^2 + a^1) 2 .\]
\[) \int ) \int\]
\end{flushleft}
\subsection*{Page 4}
\begin{flushleft}
\[\geq\]
L\par
\[C^1\]
\[C^2\]
L P\par
The diagram above shows two non-overlapping circles C1 and C2 of different sizes. The lines\par
\[\geq\]
\[L and L are the two common tangents to C^1 and C^2 such that the two circles lie on the\]
\[\geq\]
same side of each of the tangents. The lines L and L intersect at the point P which is called\par
\[the focus of C^1 and C^2.\]
\textbf{(i)} Let x1 and x2 be the position vectors of the centres of C1 and C2, respectively. Show\par
that the position vector of P is\par
\[r1x2 - r2x1\]
,\par
\[r^1 - r^2\]
\[where r^1 and r^2 are the radii of C^1 and C^2, respectively.\]
\textbf{(ii)} The circle C3 does not overlap either C1 or C2 and its radius, r3, satisfies r1 \ne r3 \ne r2.\par
\[The focus of C^1 and C^3 is Q, and the focus of C^2 and C^3 is R. Show that P , Q and R\]
lie on the same straight line.\par
\textbf{(iii)} Find a condition on r1, r2 and r3 for Q to lie half-way between P and R.\par
\end{flushleft}
\subsection*{Page 5}
\begin{flushleft}
\bigskip\textbf{Section B: Mechanics}\par
\bigskip\textbf{Question 9.} An equilateral triangle ABC is made of three light rods each of length a. It is free to rotate in\par
a vertical plane about a horizontal axis through A. Particles of mass 3m and 5m are attached\par
to B and C respectively. Initially, the system hangs in equilibrium with BC below A.\par
\textbf{(i)} Show that, initially, the angle θ that BC makes with the horizontal is given by sin θ = .\par
\textbf{(ii)} The triangle receives an impulse that imparts a speed v to the particle B. Find the\par
\[minimum speed v^0 such that the system will perform complete rotations if v > v^0.\]
\bigskip\textbf{Question 10.} A particle of mass m is pulled along the floor of a room in a straight line by a light string\par
which is pulled at constant speed V through a hole in the ceiling. The floor is smooth and\par
horizontal, and the height of the room is h. Find, in terms of V and θ, the speed of the\par
particle when the string makes an angle of θ with the vertical (and the particle is still in\par
contact with the floor). Find also the acceleration, in terms of V , h and θ.\par
Find the tension in the string and hence show that the particle will leave the floor when\par
V\par
\[tan θ = .\]
gh\par
\end{flushleft}
\subsection*{Page 6}
\begin{flushleft}
\bigskip\textbf{Question 11.} Three particles, A, B and C, each of mass m, lie on a smooth horizontal table. Particles A\par
and C are attached to the two ends of a light inextensible string of length 2a and particle B\par
is attached to the midpoint of the string. Initially, A, B and C are at rest at points (0, a),\par
(0, 0) and (0, -a), respectively.\par
An impulse is delivered to B, imparting to it a speed u in the positive x direction. The string\par
remains taut throughout the subsequent motion.\par
y\par
A\par
θ B\par
x\par
C\par
\textbf{(i)} At time t, the angle between the x-axis and the string joining A and B is θ, as shown\par
in the diagram, and B is at (x, 0). Write down the coordinates of A in terms of x, a\par
and θ at this time. Given that the velocity of B is (v, 0), show that the velocity of A\par
is (v + aθ˙ sin θ , aθ˙ cos θ), where the dot denotes differentiation with respect to time.\par
\textbf{(ii)} Show that, before A and C first collide,\par
u\par
\[3v + 2aθ˙ sin θ = u and θ˙2 = .\]
\[a^2(3 - 2 sin 2 θ)\]
\textbf{(iii)} When A and C collide, the collision is elastic (no energy is lost). At what value of θ\par
does the second collision between particles A and C occur? (You should justify your\par
answer.)\par
\textbf{(iv)} When v = 0, what are the possible values of θ? Is v = 0 whenever θ takes these values?\par
\end{flushleft}
\subsection*{Page 7}
\begin{flushleft}
\bigskip\textbf{Section C: Probability and Statistics}\par
12 Four players A, B, C and D play a coin-tossing game with a fair coin. Each player chooses\par
a sequence of heads and tails, as follows:\par
Player A: HHT; Player B: THH; Player C: TTH; Player D: HTT.\par
The coin is then tossed until one of these sequences occurs, in which case the corresponding\par
player is the winner.\par
\textbf{(i)} Show that, if only A and B play, then A has a probability of of winning.\par
\textbf{(ii)} If all four players play together, find the probabilities of each one winning.\par
\textbf{(iii)} Only B and C play. What is the probability of C winning if the first two tosses are TT?\par
Let the probabilities of C winning if the first two tosses are HT, TH and HH be p, q\par
1 1\par
\[and r, respectively. Show that p = + q.\]
2 2\par
Find the probability that C wins.\par
\bigskip\textbf{Question 13.} The maximum height X of flood water each year on a certain river is a random variable with\par
probability density function f given by\par
\[λe -λx for x \leq 0 ,\]
\[f(x) =\]
(0 otherwise,\par
where λ is a positive constant.\par
It costs ky pounds each year to prepare for flood water of height y or less, where k is a\par
\[positive constant and y \leq 0. If X \geq y no further costs are incurred but if X > y the\]
additional cost of flood damage is a(X - y) pounds where a is a positive constant.\par
\textbf{(i)} Let C be the total cost of dealing with the floods in the year. Show that the expectation\par
of C is given by\par
a\par
\[E(C) = ky + e\]
-λy\par
.\par
λ\par
How should y be chosen in order to minimise E(C), in the different cases that arise\par
according to the value of a/k?\par
\textbf{(ii)} Find the variance of C, and show that the more that is spent on preparing for flood\par
water in advance the smaller this variance.\par
\end{flushleft}
\end{document}
